\chapter{Introdução}\label{cap:introducao}

\section{Contextualização}

O projeto de placas de circuito impresso (\emph{printed circuit board} --- PCB)
ocupa posição central no desenvolvimento de produtos eletrônicos. O processo
parte da especificação funcional, passa pela captura esquemática, pela escolha
de componentes e \emph{footprints}, pela definição da pilha de camadas e pelo
roteamento das trilhas, e culmina na geração dos arquivos de fabricação ---
tipicamente Gerber, arquivo de furação e lista de materiais
\citep{kicaddocs,horowitz2015}. Cada etapa envolve decisões interdependentes:
um ajuste de geometria pode alterar a integridade do sinal, o custo de
montagem ou a viabilidade de manufatura.

Nas últimas décadas, o setor consolidou ferramentas de automação de projeto
eletrônico (EDA) que oferecem tanto interfaces gráficas quanto mecanismos de
automação --- \emph{scripts}, interfaces de programação e utilitários de linha
de comando. Esse conjunto de mecanismos é o que permite integrar o projeto a
fluxos automatizados: verificação de regras, geração de documentação e
preparação de fabricação sem intervenção manual repetitiva.

Paralelamente, assistentes baseados em grandes modelos de linguagem (LLMs)
passaram a participar de tarefas de engenharia. Tais modelos demonstram
capacidade de encadear raciocínio e ação \citep{yao2023react}, aprender a
invocar ferramentas externas \citep{schick2023toolformer} e operar interfaces
projetadas para agentes \citep{yang2024sweagent}. Quando o objeto de trabalho
é um \emph{software} de engenharia instalado localmente, a pergunta deixa de
ser apenas ``o modelo consegue raciocinar sobre o problema?'' e passa a ser
``como expor as operações da ferramenta ao modelo de forma padronizada, segura
e verificável?''.

Em novembro de 2024, a Anthropic publicou o \emph{Model Context Protocol}
(MCP), um protocolo aberto para conectar assistentes de IA a ferramentas e
fontes de dados \citep{anthropic2024mcp}. Baseado em JSON-RPC 2.0, o MCP
define papéis de anfitrião (\emph{host}), cliente e servidor, além de três
primitivas principais: \emph{tools} (operações executáveis), \emph{resources}
(dados de contexto) e \emph{prompts} (modelos de interação)
\citep{mcp2025spec,mcp2025docs}. A adoção do protocolo por diferentes
fornecedores de modelos e editores de código reduziu o custo de integração:
um único servidor pode atender diferentes assistentes.

\section{Problema e motivação}

O KiCad, suíte de EDA de código aberto, oferece três caminhos de automação: a
API Python \texttt{pcbnew} exposta via SWIG, a API IPC para comunicação com a
interface gráfica e o utilitário \texttt{kicad-cli} para verificações e
exportações \citep{swig,kicaddev,kicadcli}. Ainda assim, integrar essas
operações a um assistente de IA exige, sem uma camada padronizada, que cada
equipe escreva e mantenha \emph{scripts} próprios, com decisões repetidas de
tratamento de erro, sessão, descoberta de comandos e documentação.

O problema endereçado neste trabalho é, portanto: \emph{como expor o fluxo de
projeto de PCBs do KiCad a assistentes de IA de forma ampla, descobrível,
robusta e verificável?} A motivação é dupla: de um lado, reduzir o esforço de
integração e o risco de operações destrutivas inconsistentes sobre arquivos de
projeto; de outro, criar uma base de integração que possa ser avaliada,
testada e estendida pela comunidade.

\section{Questões de pesquisa}

O trabalho busca responder às seguintes questões:

\begin{enumerate}[leftmargin=1.4em,itemsep=0.2em]
  \item Q1 --- Como estruturar um servidor MCP que cubra o fluxo completo de
        projeto de PCB no KiCad, desde a criação do projeto até a exportação
        de fabricação?
  \item Q2 --- Como organizar e tornar descobríveis centenas de operações para
        um agente de IA, sem ocultar esquemas e sem sobrecarregar o cliente?
  \item Q3 --- Em que medida a implementação resiste a falhas previsíveis de
        integração (respostas fora de ordem, tempos-limite, estados de sessão)
        e como isso pode ser verificado de forma automatizada?
  \item Q4 --- Que evidências de qualidade e de uso real podem ser obtidas a
        partir de um projeto de placa desenvolvido no mesmo ambiente?
\end{enumerate}

\section{Objetivos}

\subsection{Objetivo geral}

Especificar, implementar e avaliar um servidor \emph{Model Context Protocol}
que exponha o fluxo de projeto de placas de circuito impresso do KiCad a
assistentes de inteligência artificial, com cobertura funcional ampla,
mecanismos de descoberta, integrações de fabricação e verificação automatizada.

\subsection{Objetivos específicos}

\begin{enumerate}[leftmargin=1.4em,itemsep=0.2em]
  \item definir a arquitetura do servidor e o protocolo interno entre as
        camadas de integração;
  \item implementar o catálogo de ferramentas MCP cobrindo esquemático,
        layout, roteamento, regras de projeto, bibliotecas, exportação e
        cadeia de suprimentos;
  \item implementar mecanismos de descoberta de ferramentas e de exposição de
        estado do projeto como recursos MCP;
  \item construir e executar suítes de teste e integração contínua que
        verifiquem o comportamento do servidor;
  \item avaliar a solução por meio de métricas de implementação e de um estudo
        de caso com uma placa real de quatro camadas;
  \item registrar limitações, pendências e direções de continuidade.
\end{enumerate}

\section{Justificativa}

A relevância do trabalho se apoia em três argumentos. Primeiro, a padronização
via MCP permite que o mesmo servidor sirva a diferentes assistentes, evitando
integrações proprietárias duplicadas. Segundo, o domínio de projeto de PCB é
rico em operações verificáveis automaticamente (regras de projeto, paridade
esquemático--placa, exportações), o que permite construir evidências
objetivas de qualidade. Terceiro, a documentação pública e os testes
automatizados reduzem a barreira de entrada para contribuições e para o uso em
contextos educacionais, como cursos de engenharia que adotam o KiCad.

\section{Delimitação e escopo}

O trabalho concentra-se no servidor MCP e no fluxo de projeto de PCB no KiCad.
Não fazem parte do escopo: o desenvolvimento de um modelo de linguagem próprio;
a avaliação de desempenho do modelo de IA em si; a medição controlada de ganho
de produtividade com usuários; e a garantia de fabricabilidade completa do
projeto usado como estudo de caso. Tais temas são retomados como limitações e
trabalhos futuros.

\section{Metodologia em síntese}

A pesquisa é aplicada, de natureza tecnológica, com desenvolvimento de sistema
e avaliação por estudo de caso instrumental. As etapas compreendem auditoria do
domínio e do código, definição de requisitos, implementação em duas camadas,
verificação por testes e integração contínua, e estudo de caso com uma placa
seguidora de linhas. O Capítulo~\ref{cap:metodologia} detalha o percurso e as
métricas adotadas.

\section{Organização do trabalho}

O Capítulo~\ref{cap:fundamentacao} apresenta a fundamentação teórica;
o Capítulo~\ref{cap:relacionados} discute trabalhos relacionados;
o Capítulo~\ref{cap:metodologia} detalha materiais e métodos;
o Capítulo~\ref{cap:arquitetura} descreve a arquitetura do servidor;
o Capítulo~\ref{cap:catalogo} apresenta o catálogo de ferramentas e os fluxos
de uso; o Capítulo~\ref{cap:estudo-de-caso} documenta o estudo de caso;
o Capítulo~\ref{cap:resultados} consolida os resultados;
o Capítulo~\ref{cap:discussao} discute os achados, limitações e ameaças à
validade; e o Capítulo~\ref{cap:conclusao} conclui o trabalho e indica
trabalhos futuros.
